% --------------------------------------------------------------------------
% File: src/architecture/decisions/adr-lifecycle-and-review-guide.tex
% Style: style v1.0 (tooling/latex/style.sty)
% Engine: pdfLaTeX through tooling/scripts/latex_build.py
% Purpose: reusable unfilled template; not an approved project decision
% --------------------------------------------------------------------------
\documentclass[11pt,letterpaper]{article}
\usepackage{technical-design-spec}
\renewcommand{\DocTitle}{ADR Governance and Review Guide}
\renewcommand{\DocSubtitle}{Lifecycle, decision rights, traceability and evidence}
\renewcommand{\DocVersion}{v1.0}
\renewcommand{\DocStatus}{Template}
\renewcommand{\DocOwner}{Architecture Governance}
\renewcommand{\DocAuthor}{Architecture Documentation Library}
\renewcommand{\DocDate}{2026-10-28}
\renewcommand{\DocSummary}{Establish a repeatable, lightweight architecture decision workflow with stable record identity, review gates, validation and supersession rules.}
\begin{document}
\makedoctocpage

\section{Purpose and authority}
An architecture decision record (ADR) captures one architecturally
significant choice, the alternatives actually evaluated, the reasoning,
consequences, accountable authority and proof of implementation. An ADR
is a decision artifact, not a general meeting note, a design catalogue or
a complete architecture description (AD). This guide states a
\emph{repository convention}, not a claim of independent certification.

\section{When to create an ADR}
Create an ADR for choices affecting externally visible interfaces,
architectural quality attributes, security boundaries, shared platform
contracts, data ownership, substantial financial or operational
commitments, or hard-to-reverse dependencies. Avoid ADR inflation for
routine code changes governed by an accepted architectural constraint.
Record the architectural decision once and link subordinate project
decisions as needed.

\section{Identity, location and authorship}
Use immutable, monotonically assigned \texttt{ADR-0001}-style IDs within
each adopting project. Store actual accepted project decisions in that
project's version-controlled \texttt{docs/adr/} directory, unless an
approved architecture-description location is already established.
This library's \texttt{templates/} directory is for blank source files
and MUST NOT be treated as a register of approved decisions.
Name project records \texttt{0001-short-decision-title.tex}, retaining
the ID inside the document. Do not recycle IDs.

\section{Roles and separation of duties}
\begin{tabularx}{\linewidth}{@{}p{0.26\linewidth}X@{}}
\toprule
\textbf{Role} & \textbf{Decision responsibility}\
\midrule
Proposer / author & Define context, non-goals, alternatives, expected consequences and evidence.\\
Accountable decision authority & Confirm scope, approve or reject; require risk-owner concurrence when necessary.\\
Architecture reviewer & Challenge assumptions, architecture interactions, concerns and consistency.\\
Security / privacy reviewer & Evaluate trust boundaries, policy obligations, risk treatment and exceptions.\\
Operations / platform reviewer & Verify supportability, recovery, telemetry, spending and exit path.\\
Implementer / verifier & Demonstrate implementation against stated acceptance criteria.\\
\bottomrule
\end{tabularx}

\section{Lifecycle and change control}
\begin{enumerate}
\item \textbf{Proposed}: assign ID, name owner, record context and options, open review and link requirements.
\item \textbf{Accepted}: record responsible approval, acceptance date, conditions and evidence source; adoption can begin.
\item \textbf{Rejected}: retain as a reviewed historical proposal where useful; document rejecting authority and rationale.
\item \textbf{Deprecated}: decision remains historically accurate but is no longer recommended for new work.
\item \textbf{Superseded}: link to replacement ADR and its approval; retain historical content for auditability.
\end{enumerate}
A decision becomes accepted through accountable review, not a status
badge or merge alone. Treat accepted record rationale as historically
stable. Correct factual mistakes with an explicit amendment and review
history; use a new ADR for materially different choices.

\section{Minimum review gate}
\begin{enumerate}
\item The decision is singular, bounded and architecturally significant;
the owner and authority are identifiable.
\item Affected stakeholders, concerns, architecture baseline and assumptions
are named with stable references.
\item Feasible alternatives (including the status quo where relevant)
and their trade-offs are evaluated using verifiable criteria.
\item Security, privacy, resilience, operability, cost and reversibility
are evaluated or explicitly justified as out of scope.
\item Acceptance criteria are measurable, mapped to implementation and
testing, and linked to evidence; open risks have owners.
\item Reviewers can trace the record to the relevant architecture views
and detect affected correspondences / known inconsistencies.
\item Supersession, revisit conditions and recovery paths are documented.
\end{enumerate}

\section{Connection to ISO/IEC/IEEE 42010:2022}
The standard governs architecture descriptions. Within an AD, retain
essential architectural decisions and rationale and maintain
cross-references to stakeholders, concerns, viewpoints, views,
view components and correspondence methods as applicable.
In this repository, the standards suite's
\texttt{05-correspondence-decisions.tex} proposes a local
\texttt{Decision(adr\_id)} logical record; the identifiers in the ADR
templates are local conventions, not standard-mandated notation.
A completed ADR supports, but does not prove, AD conformance.

\section{Links to existing governance artifacts}
\begin{itemize}
\item CCoE \emph{Decision Log 3.4}: portfolio/project decision register;
link to a detailed ADR by stable ID, rather than replacing the log.
\item CCoE \emph{Change Request 3.2}, \emph{Change Log 3.3} and
\emph{Change Management Plan 2.39}: control approved changes to delivery
or scope triggered by an ADR.
\item Cloud engineering suite \emph{T05}: concise cloud
decision worksheet; use the central cloud ADR for a fully evidenced choice.
\item ISO/IEC/IEEE 42010 architecture-description templates: maintain
the AD baseline and views; do not duplicate entire diagrams inside an ADR.
\end{itemize}

\section{Example quality indicators}
Review lead time, undocumented exceptions, unresolved impacted views,
time to verify acceptance criteria, ratio of accepted decisions with
reproducible evidence, and superseded decisions lacking bidirectional
links. Metrics support governance improvement, not automatic acceptance.

\section{Sources and normative caution}
\begin{itemize}
\item ISO/IEC/IEEE 42010:2022, \emph{Software, systems and enterprise---Architecture description}.
\item Repository standard interpretation: \\ \texttt{src/architecture/standards/iso-iec-ieee-42010/}.
\item Repository style API: \\ \texttt{tooling/styles/latex/technical-design-spec.sty}.
\end{itemize}

\end{document}
